\subsection{1D Jiang-Shu Density Pulse Advection}
\label{SEC:Testcases:JiangShu}

\begin{figure}[h!tb]
  \begin{centering}
    \begin{subfigure}{3.0in}
      \insertimage[width=\textwidth]{test_cases/jiangshu_start}
      \caption{Initial condition.}
      \label{FIG:JiangShuStart}
    \end{subfigure}%
    \begin{subfigure}{3.0in}
      \insertimage[width=\textwidth]{test_cases/jiangshu_dx_0.20}
      \caption{\PO[3] result after advecting one pulse width.}
      %%% One pulse width = $\Delta x=0.2$ m.
      \label{FIG:JiangShuDx0.20}
    \end{subfigure}
    \caption{Jiang-Shu density pulse test case.}
    \label{FIG:JiangShu}
  \end{centering}
\end{figure}

This test case features the inviscid advection of a density pulse, and can be
used to assess the numerical dissipation of the scheme.
%
\Fref{FIG:JiangShuStart} shows the starting density profile and the
domain length.
%
Both the reference conditions and the pulse shapes are preset
within the code.
%
Gauss, square, triangular, and elliptical pulse shapes with an amplitude of
1.0~\unitDens{} are superimposed onto the reference flow conditions.
%
These pulse shapes differ slightly from those described in
\sref{SEC:Testcases:DensityPulse}.
%
Periodic boundary conditions are applied.
%
\Fref{FIG:JiangShuDx0.20} shows a sample \PO[3] result using the resolution
marker to identify troubled cells and applying a filter with
$\inp[nmlLimiter::LocalFilterOrder]{Local\_Filter\_Order}=4$ to damp
oscillations.

\SeeAlso{
\inp[nmlInitialization::Initialization]{Initialization}=27,
\inp[nmlGrid::InternalGridDoF]{Internal\_Grid\_DoF},
\inp[nmlOutputVariables::ExactDensity]{Exact\_Density},
\inp[nmlOutputVariables::ExactPressure]{Exact\_Pressure}
}

